#import "@local/bootlin:0.1.0": *

#import "@local/bootlin-yocto:0.1.0": *

#import "@local/bootlin-utils:0.1.0": *

#import "../../typst/local/themeBootlin.typ": *

#import "../../typst/local/common.typ": *

#show: bootlin-theme.with( aspect-ratio: "16-9", config-common(handout: "handout" in sys.inputs and sys.inputs.handout == "1", ))

#show raw.where(block: true): set block(fill: luma(240), inset: 1em, radius: 0.5em, width: 100\%)

#show raw.where(block: true): set text(size: 11pt)

#show raw.where(block: false): r => text(fill: color-link)[#r]

#show raw.where(lang: "c", block: true): set block(fill: luma(240), inset: 0.4em, radius: 0.5em, width: 95\%, breakable: true, above: 12pt, below: 12pt)

#show raw.where(lang: "c", block: true): set text(11pt)

#show raw.where(lang: "console", block: true):set block(fill: luma(240), inset: 0.4em, radius: 0.5em, width: 95\%, breakable: true, above: 6pt)

#show raw.where(lang:"console", block: true): set text(12pt)

 #setuplabframe([U-Boot / TF-A],[
  Time to start the practical lab!
  \begin{itemize}
  \item Communicate with the board using a serial console
  \item Configure, build and install the bootloader stages:
        \begin{itemize}
		\item {\em TF-A} and {\em U-Boot} on STM32MP1 and Beagleplay
		\item {\em U-Boot SPL} and {\em U-Boot} on BeagleBoneBlack and QEMU
	\end{itemize}
  \item Learn {\em U-Boot} commands
  \item Set up {\em TFTP} communication with the host
  \end{itemize}
])
